\documentclass{report}
\usepackage{amsmath,amssymb}
\setlength{\parindent}{0mm}
\setlength{\parskip}{1em}
\begin{document}
\begin{center}
\rule{6in}{1pt} \
{\large Mahadevan Ganesh \\
{\bf An efficient preconditioned iterative algorithm for Helmhotz wave propagation models}}

Department of Applied Mathematics \& Statistics \\ Colorado School of Mines \\ Golden \\ CO 80401
\\
{\tt mganesh@mines.edu}\\
Charles Morgenstern\end{center}

We consider the Helmholtz acoustic wave propagation model in a bounded media
with an inhomogeneous impedance boundary condition on its boundary. It is
well known that the standard Galerkin variational formulation of the
Helmholtz partial differential equation (PDE) is indefinite for large
wavenumbers, and hence the Helmholtz PDE in some literature is known as
sign-indefinite.
The lack of coercivity (indefiniteness) in the standard Galerkin problem and
iterative algorithms associated finite element method (FEM) models is one
of the major difficulties for simulating wave propagation models using
iterative methods.

A recent theoretical variational formulation of the Helmholtz PDE shows
that coercivity property can be obtained using a non-standard approach.
However, the authors of the theoretical article also questioned the
practical use of their sign-definite formulation even for the constant
coefficient Helmholtz equation with impedance boundary condition. In
particular, the key open issue is: Whether the recent sign-definite
formulation can alleviate some of the associated iterative algorithm.

Through various computer simulations, we provide a concrete answer that the
theoretically analyzed sign-definite formulation does not alleviate the key difficulty
of reducing the GMRES iterations of the associated FEM model.
We subsequently develop a new class of efficient preconditioned iterative
Galerkin FEM wave propagation models for the Helmholtz equation.


\end{document}
